\documentclass[11pt]{article}
\usepackage[margin=2.2cm]{geometry}
\usepackage{amsmath,amssymb,algorithm,algpseudocode,listings}
\usepackage{booktabs}
\usepackage{xurl}
\usepackage{xcolor}
\usepackage[colorlinks=true,linkcolor=blue!50!black,citecolor=blue!50!black,urlcolor=blue!50!black]{hyperref}
\usepackage{orcidlink}

\lstset{
    language=Python,
    basicstyle=\ttfamily\small,
    keywordstyle=\color{blue!60!black}\bfseries,
    stringstyle=\color{green!45!black},
    commentstyle=\color{gray}\itshape,
    breaklines=true,
    frame=single,
    framerule=0.3pt,
    xleftmargin=1em,
    columns=fullflexible,
    keepspaces=true,
    showstringspaces=false
}
\lstdefinestyle{appendix}{
    basicstyle=\ttfamily\tiny,
    numbers=left,
    numberstyle=\tiny\color{gray},
    stepnumber=1
}

\title{RosettaMath: Reading, Running and Proving Mathematics from \LaTeX{}\\
\large A self-hosting translator, an interactive explorer, and a dependent-type
micro-kernel, as a foundation for open mathematics education from school to research}
\author{B.S. Hartshorn \orcidlink{0009-0004-2853-655X}\\
\small \url{https://github.com/brentharts/RosettaMath}}
\date{September 2026}

\begin{document}
\maketitle

\begin{abstract}
Mathematics is written in \LaTeX{} and computed in Python, and the two are
connected by hand. RosettaMath closes that gap with three small, dependency-free
Python programs that treat \LaTeX{} as an \emph{input} language rather than a
typesetting target. \texttt{rosettamath.py} translates a strict subset of
\LaTeX{}---math mode and \texttt{algpseudocode}---into executable Python and back
again; it is self-hosting, written in the very subset it translates, and
bootstraps to a byte-identical fixed point, while the reverse direction is
verified by round-tripping functions and checking that they still behave the
same. \texttt{rosettaui.py} renders an equation glyph by glyph in a desktop
window in which every symbol is a live object: hover it for its name, origin and
conventional meaning, click it for an offline explanation, right-click it for
the encyclopedia article; the whole equation is identified against a library of
well-known physics and mathematics, and the Python it becomes is shown alongside
with physical constants resolved to \texttt{scipy.constants}.
\texttt{lean4.py} is a Calculus of Constructions kernel with de Bruijn
indices, inductive families with generated recursors---equality among
them---an untrusted
elaborator that infers implicit arguments by Miller pattern unification with
postponed constraints, and a \LaTeX{} statement language, so that a theorem
about a Python function is stated as
\verb|\forall x \in \text{Nat}, x = x|, or proved by induction, and checked
by a kernel small enough to read in one sitting. We argue that these three pieces, precisely
because they are small, local, open and written in the language students
already learn, form a practical foundation for mathematics and programming
education at every level---from a child clicking on a Greek letter to an
undergraduate proving symmetry of equality to a researcher graduating to
Lean---and that the same small trusted kernel points toward verified systems
software when its statement language becomes the contract language of a
self-contained compiler toolchain. We situate the work against the rapid
mechanisation of mathematics, in which AI systems now formalise Fermat's Last
Theorem in eleven days, and argue that what that moment most needs is not
another prover but tools with which people can read, run and check
mathematics on their own machines.
\end{abstract}

\section{Introduction}

The universal language of theoretical research is mathematics, and the
universal medium for writing it down is \LaTeX{}. The universal language for
investigating it computationally is Python. Between the two sits a manual
translation step that is slow, error-prone, and---more importantly for
education---invisible: the reader of a paper sees the equation, the user of a
simulation sees the code, and nobody sees the correspondence. The notation is
dense, overloaded and dependent on unstated convention; the program is verbose,
explicit and dependent on names. A student who has learned one has learned
almost nothing about the other.

RosettaMath began as a small attempt to make that correspondence mechanical.
Its first component, \texttt{rosettamath.py}, translates a strict subset of
\LaTeX{} into Python with no dependencies, and demonstrates the adequacy of the
subset by being written in it: a plain-Python stage~0 translates the \LaTeX{}
source of the translator into a stage~1, and stage~1 translates the same source
again, byte for byte, to a fixed point. The earlier draft of this paper
described that translator alone. Since then the project has grown two further
components, and with them a claim.

The claim is about education, and it is made at a particular moment. In the
past eighteen months the mechanisation of mathematics has moved faster than
almost anyone predicted. AlphaProof reached silver-medal performance at the
International Mathematical Olympiad with proofs checked in Lean
\cite{hubert2025}; a succession of open models followed
\cite{ren2025,wang2025kimina,chen2025seed,lin2025goedel,achim2025aristotle};
the strong Prime Number Theorem, which had stalled for eighteen months of expert
effort, was formalised by an autoformalisation agent in three weeks
\cite{mathinc2025gauss}; AI systems began to find counterexamples to conjectures
that had stood for a century \cite{buzzard2026counter}; and days before this
paper was finished, a swarm of language-model agents produced the first complete,
computer-checked formalisation of Fermat's Last Theorem in eleven days---thirteen
million lines of Lean, checked by the Lean kernel and then again by an independent
kernel implementation \cite{anthropic2026flt,newscientist2026flt}. The
mathematical community is asking, with some urgency, what its values, practice
and teaching should look like in such a world \cite{commelin2026,sloman2026,simons2026}.

We take a narrow position within that conversation. The provers are being built,
at enormous scale and cost, by organisations that can afford to. What is not
being built at the same pace is the other side: tools with which a person, on
their own computer, with no account and no server, can \emph{read} an equation
and understand its conventions, \emph{run} it as code and see the numbers, and
\emph{check} a claim about that code with a kernel small enough that they could
read the kernel itself. Those are the three components of RosettaMath, and the
argument of this paper is that they are the right foundation for teaching
mathematics and programming together, at every level, in an age when the
proving will increasingly be done by machines and the understanding must still
be done by people.

\medskip
\noindent The contributions since the previous draft are:
\begin{itemize}
\item a \textbf{self-hosting translator} extended with implicit multiplication,
inferred parameters, and a reverse translator whose output is verified by
behaviour rather than by eye (\S\ref{sec:rosettamath});
\item an \textbf{interactive explorer}, \texttt{rosettaui.py}, in which every
glyph of a typeset equation is a clickable object backed by an offline
knowledge base of 134 symbols, 44 equations and 17 concept articles, with
automatic identification of what an equation ``looks like'' (\S\ref{sec:rosettaui});
\item a \textbf{dependent-type micro-kernel}, \texttt{lean4.py}, that checks
proofs about Python functions against theorems stated in \LaTeX{}, with an
elaborator separated from a trusted core, implicit arguments, pattern
unification, constraint postponement, and inductive families with generated
recursors, so that $0 + n = n$ is proved by induction and equality itself is a
declared type rather than an axiom (\S\ref{sec:lean4});
\item an \textbf{argument} that these pieces, being local, open and modifiable,
are a better foundation for mathematics education than the client--server
applications and proprietary systems that currently dominate it
(\S\ref{sec:education}), and a sketch of how the same kernel extends to verified
systems software (\S\ref{sec:beyond}).
\end{itemize}

\section{Background}

\subsection{Literate programming and the notation gap}

Knuth's literate programming \cite{knuth} proposed that a program and its
explanation be one document. RosettaMath takes the proposal one step further
in the specific case of mathematics: the equation \emph{is} the program. Prior
systems approach this from the computational side---SymPy \cite{meurer} emits
\LaTeX{} from expressions it already holds, and PythonTeX \cite{poore} lets
\LaTeX{} documents execute embedded Python---so in both the source of truth is
code. RosettaMath reverses the direction: the source of truth is the notation a
mathematician actually writes, and the code is derived from it.

\subsection{Lean, and the bridge from proofs to prose}

Lean \cite{demoura2021} is a proof assistant and programming language built on
the Calculus of Constructions \cite{coquand1988}, whose mathematical library
Mathlib now runs to millions of lines. The Xena project's early work on
rendering Lean proofs into \LaTeX{} and HTML \cite{massot2019} recognised that
formal proofs are unreadable to mathematicians unless something translates
them back into the language mathematicians use; Massot's later Verbose Lean
\cite{massot2024} takes the pedagogical version of that idea to
undergraduates, letting students write proofs in a controlled natural language
that transfers directly to paper, with a teacher-adjustable level of
assistance. Bottoni, Cattaneo and Sacikara report a course teaching foundations
of mathematics with Lean and observe its influence on students' understanding
of proof \cite{bottoni2025}.

RosettaMath's \texttt{lean4.py} flips the Xena bridge, as its docstring says:
\LaTeX{} is the input language, and the thing being proved is a Python
function. It is not a competitor to Lean and does not attempt to be; it is a
kernel of the same design, small enough to be a teaching object, positioned
before Lean on the path a student walks.

\subsection{Machines that prove, and what they change}

The systems cited in the introduction share a structure: a language model
proposes, and a Lean kernel disposes. That structure matters for education in
two ways. First, it makes a small trusted kernel the load-bearing component of
the entire enterprise---de Moura calls this the trust bottleneck
\cite{peterman2024}, and the Fermat formalisation was checked not once but by
two independent kernels, precisely because the kernel is the only thing anyone
has to believe \cite{anthropic2026flt}. A student who has read a kernel understands
what a machine-checked proof \emph{is} in a way that a student who has only
used one does not. Second, it changes what mathematicians are for. Buzzard,
reflecting on a summer in which AI agents disproved a conjecture of Erd\H{o}s
and then the Jacobian conjecture, describes a discipline being
``outcounterexampled'' and a community whose trust in human-checked argument was
already fragile \cite{buzzard2026counter}. Commelin, Jamnik, Ochigame, Taelman
and Venkatesh, writing for the Notices of the AMS, call for the community to
safeguard its intellectual autonomy, broaden curricula, and build
academically-oriented infrastructure rather than let the future happen to it
\cite{commelin2026}; SorryDB \cite{letson2026} is one such piece of
infrastructure, a benchmark drawn from real formalisation projects rather than
competitions. The Simons Foundation's account of Lean's growth notes that the
Lean Focused Research Organization intends to transition into a foundation in
the footsteps of those behind Rust and Linux \cite{simons2026}. The
infrastructure of verified mathematics is converging on the governance model of
open-source systems software. We think its educational tools should too.

\subsection{Education, and the design of the tool}

Wegerif's recent essay on researching the future of education
\cite{wegerif2026} makes a point that is central to our argument. A widely-cited
randomised trial of nearly a thousand secondary mathematics students
\cite{bastani2025} found that those given unrestricted access to a GPT-4
assistant performed better while using it and substantially worse once it was
removed; but the same trial found that a more carefully constrained tutor
performed as well as the control. Wegerif's reading is that the design of the
tool, not the presence of AI, is what determines whether students learn---and
that an experiment which measures learning only by the old examination cannot
see the new capabilities a new tool might teach. He offers an analogy: students
taught with calculators who are then examined on an abacus will do worse on
the abacus, and this tells us nothing about whether they learned mathematics.

A proof checker is the most carefully constrained tutor there is. It will not
finish the student's proof; it will only say whether the proof the student
wrote is right, and why not. A translator that shows the Python an equation
becomes does not do the student's thinking; it shows the student what their
notation means. These are tools in the constrained category, and the argument
of \S\ref{sec:education} is that they teach.

\section{\texttt{rosettamath.py}: the self-hosting translator}
\label{sec:rosettamath}

\subsection{The subset}

The translator accepts \LaTeX{} math mode and the \texttt{algpseudocode}
environment. Table~\ref{tab:subset} summarises the correspondence. Everything
in the left column is ordinary \LaTeX{} that typesets as one would expect; the
translation is a second reading of the same text.

\begin{table}[h]
\centering\footnotesize
\begin{tabular}{@{}lll@{}}
\toprule
\LaTeX{} & Python & Note \\
\midrule
\verb|\gets| & \verb|=| & assignment; \verb|=| itself becomes \verb|==| \\
\verb|\land  \lor  \lnot  \in  \notin| & \verb|and or not in| & with their own spacing \\
\verb|\frac{a}{b}| & \verb|(a)/(b)| & parenthesised: the stacking was the grouping \\
\verb|a^{b}| & \verb|a**(b)| & \\
\verb|2x  2\pi| & \verb|2*x  2*pi| & a number touching a name \\
\verb|k_B c  m \hbar| & \verb|k_B*c  m*hbar| & a space between two atoms (new) \\
\verb|\texttt{..}| & \verb|'..'| & with \verb|\textbackslash \{ \_ \$| unescaped \\
\verb|\rho| & \verb|rho| & resolved in the caller's scope \\
\verb|\State $x \gets 1$| & \verb|x = 1| & one statement per line \\
\verb|\If \ElsIf \Else \EndIf| & \verb|if elif else| & indentation from nesting depth \\
\verb|\For{$i = a$ to $b$}| & \verb|for i in range(a, b + 1)| & inclusive, as in mathematics \\
\verb|\While  \Function  \Return| & \verb|while def return| & \\
\verb|$f(x) = 2x + 1$| & \verb|def f(x): return 2*x + 1| & an equation is a function \\
\verb|$S = k_B \log W$| & \verb|def S(k_B, W): ...| & a bare left side infers its arguments (new) \\
\verb|\begin{cases} .. \end{cases}| & \verb|if ..: return ..| & piecewise is a branch \\
\bottomrule
\end{tabular}
\caption{The translated subset. The two rows marked new are the semantic
additions since the previous draft.}
\label{tab:subset}
\end{table}

\subsection{Two semantic additions}

\paragraph{Implicit multiplication.} Juxtaposition means ``multiply'' to a reader
and nothing at all to a parser. The earlier translator handled only a numeral
touching a name; \verb|k_B c^3 A| became the syntactically invalid
\verb|k_B c**3 A|. A new predicate decides whether a literal space is a product:
the character before it must be alphanumeric, an underscore, or a closing
bracket, and what follows must begin a name, a digit, or a command whose
translation begins with a letter. The last condition is what keeps it safe.
\verb|\cdot| translates to \verb|*|, \verb|\gets| to \verb|=|, and \verb|\land|
to \verb|and| with its own surrounding spaces, so none of them is mistaken for
an operand. The test that matters is that the translator's reading of its own
\LaTeX{} source is byte-for-byte unchanged by the rule.

\paragraph{Inferred parameters.} An equation whose left-hand side is a bare
symbol---$S = \frac{k_B c^3 A}{4 G \hbar}$, written as
\verb|$S = \frac{k_B c^3 A}{4 G \hbar}$|---names no arguments, and
previously produced the invalid \verb|def S:|. The translator now builds the
function body first and then takes the free names of that body, in order of
first appearance, as the parameter list---skipping string literals, attribute
access, anything followed by an opening parenthesis, and Python keywords.
Names already present in the caller's scope stay constants:
\begin{lstlisting}
latex2py(r'$S = \frac{k_B c^3 A}{4 G \hbar}$')
# def S(k_B, c, A, G, hbar): return (k_B*c**3*A)/(4*G*hbar)
latex2py(r'$S = \frac{k_B c^3 A}{4 G \hbar}$',
         {'k_B': k, 'c': c, 'G': G, 'hbar': hbar})
# def S(A): return (k_B*c**3*A)/(4*G*hbar)
\end{lstlisting}
The second form is the point of the exercise. Fed CODATA values it gives
$1.449\times10^{54}$\,J/K for the entropy of a one-solar-mass black hole,
against a literature value of about $1.5\times10^{54}$. The equation as
published is the function that runs.

\subsection{Self-hosting, and what it costs}

The translator exists twice: in plain Python, and in the \LaTeX{} subset it
translates. Stage~0 reads the \LaTeX{} and produces stage~1; stage~1 reads the
same \LaTeX{} and must produce itself. The fixed point is currently 285 lines of
generated Python, and it is reached on every run of the test suite.

The practical consequence, which the earlier draft did not spell out, is that
every change to the translator has to be made twice---once in the Python of
stage~0 and once in the \LaTeX{} of the self-hosted source---and the fixed point
is the proof that the two agree. Both semantic additions above were made that
way. This is a cost, and it is the right cost: a translator that can only
express what it can translate is a translator whose subset is honest. The
hand-written \LaTeX{} source, typeset, appears in the appendix; a reader can
trace a \verb|while| loop from pseudocode to Python by looking left and right.

\subsection{The reverse direction, verified by behaviour}

\texttt{Python2Tex} walks a Python abstract syntax tree back into
\texttt{algpseudocode}. The previous draft described it as complete; it was
not, and the failure mode is instructive. \LaTeX{} that reads correctly can
mean something else. Precedence was ignored, so \verb|(a + b) * c| came out as
$a + b \cdot c$, a different function. Attribute access was typeset in
\verb|\texttt|, which the forward translator reads as a string literal, so
\verb|self.result| came back as \verb|'self.result'|. A chained assignment
\verb|a = b = 1| became tuple unpacking. None of these produced a \LaTeX{}
error; all of them produced wrong code.

The rewritten translator makes precedence explicit with parentheses, and
handles a subtler problem: inside \verb|$...$| a space between two names is now
implicit multiplication, so \verb|x is None| would read back as
\verb|x*is*None|. Rather than change the forward translator, the reverse one
steps outside math mode for Python keywords---\verb|$x$ is $None$|---which is
the convention the hand-written \LaTeX{} source already uses. Underscores are
escaped, since \verb|__name__| in math mode is a double subscript that \LaTeX{}
rejects; a standalone \verb|\Comment| is hung off a \verb|\State|, since
\texttt{algpseudocode} treats it as an attachment. Both of those were found by
compiling the output, which nothing had done before.

The test is the deliverable. Twenty-one functions are translated to \LaTeX{},
read back with the forward translator, executed, and compared against the
originals on sample arguments---by result, not by text, because \verb|\frac|
legitimately comes back with parentheses it did not have. Every function in
\texttt{rosettamath.py} survives the journey, including the forward translator
itself, which still behaves identically after being written out as pseudocode
and read back in. Constructs the subset cannot express---\verb|with|,
\verb|try|, decorators---are reported as a comment and recorded in a warnings
list, never dropped in silence.

\section{\texttt{rosettaui.py}: every glyph is an object}
\label{sec:rosettaui}

\subsection{What it does}

The explorer renders an equation glyph by glyph into a desktop window
(figure~\ref{fig:ui}). Each glyph is a live object. Hovering it colours it and
shows a tooltip with the symbol's name, its Greek or Latin origin, and what it
conventionally denotes in physics and mathematics; a left click opens a fuller
offline description; a right click offers the relevant encyclopedia article,
the \LaTeX{} to copy, and a submenu of other equations in the library that use
the same symbol. Roles are colour-coded---Greek letters, constants, relations,
variables---which makes visible a distinction that typesetting leaves implicit
and that newcomers find hardest.

\begin{figure}[h]
\centering
\fbox{\parbox{0.92\textwidth}{\small
\textbf{Rendered:}\quad
$\displaystyle S = \frac{k_B\, c^3 A}{4\, G\, \hbar}$
\hfill
\textbf{This equation:}\; looks like the Bekenstein--Hawking entropy (General relativity).
A black hole's entropy is proportional to the area of its horizon, not its volume\ldots\\[3pt]
\textbf{Python (rosettamath):}\;
\texttt{from scipy.constants import k as k\_B, c, G, hbar}\;
\texttt{def S(A): return (k\_B*c**3*A)/(4*G*hbar)}
}}
\caption{The three panels of the explorer, for one equation. In the running
program the rendered equation is interactive: $\hbar$ is purple because it is a
constant, $=$ is brown because it is a relation, and hovering any of them
explains it.}
\label{fig:ui}
\end{figure}

The whole equation is classified. A feature set is extracted from the parsed
tree---the symbols present, plus structural tags such as ``has a fraction'',
``has a Laplacian'', ``has an exponent of two''---and matched against a library
of 44 equations, each with must-have and confidence-raising features. The
result is a verdict such as ``this looks like the Klein--Gordon equation'', or a
family fallback such as ``a second-order partial differential equation of the
kind that governs fields---compare the wave, heat and Laplace equations, which
differ only in what the Laplacian is set equal to''. The classifier identifies
both forms of the Schr\"odinger equation, Klein--Gordon, Dirac, the Lorentz
factor, the Debye length, and the d'Alembert operator, among others; the
library spans classical mechanics, electromagnetism, quantum mechanics,
relativity, plasma physics, thermodynamics, information theory and probability.

A third panel shows what \texttt{rosettamath.py} makes of the same \LaTeX{}.
Symbols that are physical constants are recognised as such, kept out of the
argument list, and bound with a generated \verb|from scipy.constants import|
line. That is the symbol-to-library mapping the earlier draft described as an
aspiration, now doing something.

\subsection{The knowledge base}

The offline knowledge base holds 134 symbols, 44 equations and 17 longer
concept articles, cross-linked and with 165 distinct encyclopedia links. Each
symbol entry gives the glyph, the \LaTeX{}, the name, a category, and a
description written to explain \emph{convention} rather than syntax: what
$\rho$ means in four different sub-fields, why the curly $\partial$ differs
from the straight $d$, why the $i$ in the Schr\"odinger equation cannot be
removed. The concept articles are short essays on the things a syntax
reference cannot tell you---overloaded notation, the three notations for a
derivative, the nabla operator, the Laplacian as a sorting principle for the
great partial differential equations, bra--ket notation, tensor index
conventions, conservation laws and symmetry, the action principle, entropy,
Fourier analysis, dimensional analysis, natural units, and the d'Alembertian.
An alias table folds alternative spellings---\verb|\square|, \verb|\dfrac|,
\verb|\ne|---onto the entries that document them, so each idea appears once in
the menus and classifies identically however it is written.

The menu bar exposes all of this as a small browsable encyclopedia: equations
by field, concepts by category, symbols by category, each with a filterable
browser. The point is not that this is a large knowledge base---it is
not---but that it is a \emph{local} one, in a single Python file, that a
teacher can open and extend with the symbols and equations their own course
uses.

\subsection{Rendering, honestly}

Rendering \LaTeX{} in a desktop widget is harder than it looks. The explorer
uses its own two-pass layout engine---measure the tree, then place real
items---so that fractions and roots size themselves around their contents. The
radical is drawn as a path that stretches, since a fixed glyph degenerates over
tall content. A detail worth reporting: Latin Modern Roman, the natural choice
of typeface, has no glyph for sixty-seven of the symbols the explorer uses,
including every Greek letter, $\nabla$, $\partial$ and $\hbar$. They rendered
anyway, because the font system silently substituted, but the metrics were
wrong. Glyphs are now measured and drawn with a typeface chosen per character,
serif for Latin and digits and Latin Modern Math for operators and Greek, so
that measuring and drawing agree.

For anything the unicode approximation cannot stack, the explorer shells out
to \texttt{pdflatex} and shows the result as an image. It tries poppler before
ImageMagick, because ImageMagick delegates PDF to Ghostscript, which is often
absent; a dependency checker reports which tools are present rather than
failing. Everything is cached by content hash.

\subsection{Why a desktop program}

The explorer is 3,171 lines of Python and PyQt5, in one file, with no server,
no account, no network requirement beyond an optional browser link, and no
build step. This is a design position, argued in \S\ref{sec:education}. Here we
note only its immediate consequence: a teacher who wants to add a symbol edits a
list; a student who wants to change the colour of Greek letters edits a
dictionary; a class that wants a new panel writes a method. The program is
small enough to be a thing one modifies rather than a thing one uses.

\section{\texttt{lean4.py}: a micro-kernel for proofs about Python}
\label{sec:lean4}

\subsection{The kernel}

\texttt{lean4.py} implements the Calculus of Constructions
\cite{coquand1988}---the foundation Lean and Coq are built on---in under two
thousand lines including its tests. The term language is universes, variables,
$\Pi$-types, $\lambda$-abstractions and application; the checker implements the
standard rules, with an impredicative \textsf{Prop} and no type-in-type.

Bound variables are de Bruijn indices \cite{debruijn1972}; free variables and
constants keep their names, and a binder stores a name only as a printing hint.
The first version of the kernel used names throughout and had the two bugs
that representation invites: substitution captured variables, so a correct
proof could be made to typecheck for the wrong reason, and $\forall x{:}\mathrm{Nat},
\mathrm{Nat}$ was judged different from $\forall y{:}\mathrm{Nat}, \mathrm{Nat}$, so
a correct proof could be rejected as soon as a statement and a function named
their variables differently. De Bruijn indices do not fix those bugs; they make
them unstateable. A name in a substituted term cannot collide with an index,
and alpha-equivalent terms are the same object, so structural equality is
alpha-equivalence with no renaming machinery. The constructors still take
names---\verb|Lambda("T", Universe(1), Var("T"))| abstracts the \verb|T| for
you---so terms read as they always did. One consequence deserves mention: the
printer renames a binder that would shadow a free variable, because a checker
whose output cannot be trusted to mean what it says is worse than none.

\subsection{Theorems about Python functions}

A Python function is compiled to a kernel term by an AST visitor: parameters
with type annotations become nested $\lambda$s, calls become curried
application, and the single \verb|return| is the body. The statement is written
in \LaTeX{}:
\begin{lstlisting}
@theorem(r'\forall x \in \text{Nat}, x = x')
def reflexivity(x: 'Nat'):
    return refl(x)
\end{lstlisting}
The decorator compiles the function, elaborates it, reads the statement as a
type, and requires the two to agree definitionally. A theorem that does not
prove what it claims raises; \verb|make proofs| gates the build. A
\verb|--non-strict| mode records the failure and continues. The earliest
version of this decorator accepted the statement and never read it, checking
only that the Python was well typed; its flagship example ``proved'' $\forall
x, x = x$ by returning $x$, whose type is $\mathrm{Nat}$, and it passed because
$\mathrm{Nat}$ had been declared a proposition. We mention this because it is
the exact failure mode the tool exists to prevent, and it was caught by making
the tool honest.

The statement language is a recursive-descent parser over the same tokeniser
the explorer uses: \verb|\to| is a right-associative function type,
\verb|\forall x \in A, B| and \verb|\forall (x : A), B| are dependent function
types, \verb|\forall \{A : \text{Type}\}, B| declares an implicit binder,
juxtaposition is application, \verb|\text{Nat}| and \verb|\mathbb{N}| name
types, and \verb|a = b| elaborates to $\mathrm{Eq}\;A\;a\;b$ with the carrier
taken from the enclosing binder or left as a hole for the elaborator.

\subsection{An elaborator, separated from the kernel}

The kernel knows nothing about implicit arguments; a hole reaching it is an
error. In front of it sits an elaborator that fills the holes and hands the
completed term to the kernel for independent verification---the same
separation Lean maintains, and the property that makes the design safe to
teach with: a bug in the elaborator can cost a confusing error message, never a
false theorem. The elaborator went through three stages, each forced by an
example.

\paragraph{Implicit arguments.} \verb|refl| has type $\forall \{A{:}\mathrm{Type}\},
\forall a{:}A, \mathrm{Eq}\;A\;a\;a$. Each use contributes a hole for $A$, solved
by unifying the expected argument type with the actual one, so the proof is
written \verb|refl(x)|, as Lean writes it, rather than \verb|refl(Nat, x)|.
Elaboration runs over \emph{opened} terms: entering a binder replaces the bound
variable with a fresh free name and the binder is closed again on the way
out. That is what lets a hole be solved with a bound variable---as in
\verb|refl(a)| where $a{:}A$ and $A$ is itself bound---since a de Bruijn index
would mean something different at every depth it appeared, and a name does
not. \verb|explicit(refl)(Nat, x)| is Lean's \verb|@refl|.

\paragraph{Pattern unification.} An implicit argument in a \emph{dependent}
position is only ever seen applied. The transport constant has the type
\[\forall \{A\}\{P{:}A\to\mathrm{Prop}\}\{a\,b{:}A\},\quad
\mathrm{Eq}\;A\;a\;b \to P\,a \to P\,b,\]
so solving for the motive $P$ means answering $?P\,a = \mathrm{Eq}\;A\;a\;c$,
a higher-order question that first-order decomposition can only answer by
demanding $a = c$. Unification therefore covers Miller's pattern fragment
\cite{miller1991}: a hole applied to distinct variables is solved by abstracting
them out of the other side. Full higher-order unification is undecidable, so
outside the fragment the elaborator refuses rather than guesses. Checking is
bidirectional, with the expected type pushed inward through the $\lambda$s, since
two closed types facing each other at the top would leave indices with no name to
abstract over. This is enough for
\begin{lstlisting}
@theorem(r'\forall \{A : \text{Type}\}, \forall a \in A, \forall b \in A, \forall c \in A, '
         r'\text{Eq} A a b \to \text{Eq} A a c \to \text{Eq} A b c')
def transitivity(A: 'Type', a: 'A', b: 'A', c: 'A',
                 h: r'\text{Eq} A a b', p: r'\text{Eq} A a c'):
    return transport(h, p)
\end{lstlisting}
where the elaborator synthesises the motive $\lambda z.\,\mathrm{Eq}\;A\;z\;c$.

\paragraph{Postponed constraints.} Even inside the fragment a constraint can
have several legal answers. $?P\,a = \mathrm{Eq}\;A\;a\;a$ admits
$\lambda z.\,\mathrm{Eq}\;A\;z\;z$, $\lambda z.\,\mathrm{Eq}\;A\;z\;a$,
$\lambda z.\,\mathrm{Eq}\;A\;a\;z$ and the constant, whenever $a$ is in scope
where the hole was made; so each hole records the scope it was created in, and
a constraint is ambiguous only when its variable is in it. Ambiguous constraints
are postponed and settled once every constraint on that hole is known. Each
proposes the whole family of readings, obtained by abstracting any subset of the
occurrences, most abstracted first; a reading is accepted only if it satisfies
all of them. That is what lets symmetry be written the natural way:
\begin{lstlisting}
@theorem(r'\forall \{A : \text{Type}\}, \forall a \in A, \forall b \in A, '
         r'\text{Eq} A a b \to \text{Eq} A b a')
def symmetry(A: 'Type', a: 'A', b: 'A', h: r'\text{Eq} A a b'):
    return transport(h, refl(a))
\end{lstlisting}
Solved eagerly, the first constraint's obvious answer, $\lambda z.\,\mathrm{Eq}\;A\;z\;z$,
wins, and the theorem fails with \emph{stated} $\mathrm{Eq}\;A\;b\;a$, \emph{proved}
$\mathrm{Eq}\;A\;b\;b$. With postponement the second constraint's proposal
$\lambda z.\,\mathrm{Eq}\;A\;z\;a$ survives both and is chosen. When several
readings survive, the theorem is still proved---the kernel checks the finished
term either way---but the elaborator reports that the choice was not forced.
A hole no reading can satisfy is reported, not left tentative.

\subsection{Definitions, inductive types, and induction}

A global name may now carry a value as well as a type, so a definition unfolds
during normalisation (delta reduction); and an inductive type may be declared
with its constructors, from which the recursor and its computation rule are
generated (iota reduction). $\mathrm{Nat}$ is declared rather than assumed:
\begin{lstlisting}
inductive(env, 'Nat', [('zero', []), ('succ', [REC])])
\end{lstlisting}
which yields \verb|zero|, \verb|succ|, a recursor \verb|Nat.rec| eliminating
into $\mathrm{Type}$ for defining functions, and an induction principle
\verb|Nat.ind| eliminating into $\mathrm{Prop}$ for proving theorems. There are
two because there is no universe polymorphism here; they share one computation
rule. Definitions may not mention themselves---recursion belongs in the
recursor---so unfolding always terminates.

A family may also take \emph{parameters}, fixed across all constructors, and
\emph{indices}, which vary from one constructor to the next. Parameters give
$\mathrm{List}$: they are implicit in the constructors, so a list is written
\verb|cons(7, cons(8, nil))| and its element type is inferred rather than
spelled out, and a \verb|length| defined from \verb|List.rec| reduces
$\mathrm{length}\;(\mathrm{cons}\;7\;(\mathrm{cons}\;8\;\mathrm{nil}))$ to $2$.

Indices are what make equality expressible rather than assumed:
\begin{lstlisting}
inductive(env, 'Eq', [('refl', [], [Var('a')])],
          params=[('A', Universe(1)), ('a', Var('A'), False)],
          indices=[('b', Var('A'))], level=0)
\end{lstlisting}
$\mathrm{Eq}\;A\;a\;b$ is the family whose single constructor only ever builds
the case where the index $b$ is $a$---which is the entire content of
equality---and whose recursor is the J rule. Consequently
\verb|symm|, \verb|trans|, \verb|congrArg| and \verb|transport|, which were
declared as axioms in the previous version of this system, are now
\emph{proved} from that recursor and checked like any other definition. The
distinction is not decorative: an axiom is something a reader has to trust, and
a definition is not. The trusted base is the kernel and nothing else.

With addition defined by recursion on its second argument, the two halves of a
familiar pair of facts come apart in a way that is worth showing a student.
$\mathrm{add}\;2\;3$ computes to $5$, and $m + 0 = m$ holds by computation
alone, so \verb|refl| proves it; the checker reports the claim as
$\mathrm{Eq}\;\mathrm{Nat}\;(\mathrm{add}\;m\;\mathrm{zero})\;m$ and what was
proved as $\mathrm{Eq}\;\mathrm{Nat}\;m\;m$, making visible that the gap was
closed by reduction. But $0 + n$ is stuck until $n$ is a constructor, so
$0 + n = n$ requires the induction principle:
\begin{lstlisting}
@definition(r'\text{Nat} \to \text{Prop}')
def add_zero_motive(k: 'Nat'):
    return Eq(Nat, add(zero, k), k)

@theorem(r'\forall n \in \text{Nat}, \text{Eq} \text{Nat} (\text{add} \text{zero} n) n')
def add_zero_left(n: 'Nat'):
    return explicit(Nat.ind)(add_zero_motive, refl(zero), add_zero_step, n)
\end{lstlisting}
That two superficially symmetric statements need different proofs, and that the
asymmetry comes from which argument the definition recurses on, is a lesson
about computation that a student can now discover by trying both.

\verb|@definition| adds a checked Python function to the environment, so a
proof can build on earlier ones; before this the environment was a fixed list
of constants and every theorem stood alone.

\subsection{What it does and does not do}

The environment ships the inductive types $\mathrm{Nat}$, $\mathrm{Bool}$,
$\mathrm{List}$ and $\mathrm{Eq}$, with \verb|refl| as a constructor and
\verb|symm|, \verb|trans|, \verb|congrArg| and \verb|transport| proved from
\verb|Eq.ind|; the test suite runs 153 checks, each a regression from a bug the
kernel used to have or a construction it can now do, and each of the
elaboration features is verified to be doing the work by disabling it and
watching the corresponding theorem fail. The kernel has no universe
polymorphism and no tactics; every proof is a term. A recursive argument is
allowed only in a family without indices, since the induction hypothesis would
otherwise have to name the indices of that occurrence. The constraint solver is
greedy across holes. None of that is hidden, and all of it is the point: a
student can read \texttt{lean4.py} in an afternoon and understand exactly what
it means for a machine to have checked a proof, and then go to Lean knowing
what the big system is doing and why.

\section{Education, from school to research}
\label{sec:education}

\subsection{A progression}

The three components correspond to three things a person learning mathematics
needs to be able to do with a piece of notation: read it, run it, and check a
claim about it. They also correspond, roughly, to a progression of age and
level, though the boundaries are deliberately porous.

\paragraph{Reading (school).} A child learning that $E = mc^2$ means something
can hover the $c$ and learn that it is the speed of light, a conversion factor
between space and time and the universal speed limit; hover the $=$ and learn
that in a programming language this same symbol means assignment and a
different symbol means comparison, which is a fact worth knowing at ten. The
explorer's tooltips are written for that reader. Greek letters have entries
that say what they \emph{are} before what they denote. The equation library is
a gallery to browse, with the guarantee that each item comes with the Python it
becomes---so the student who asks ``but what does it \emph{do}?'' gets an
answer in a language with a \verb|for| loop in it.

\paragraph{Running (secondary and undergraduate).} The translator's rules are
themselves a curriculum. That $\frac{a}{b}$ must become \verb|(a)/(b)| teaches
that the visual stacking was doing the grouping. That a sum from $1$ to $n$
becomes \verb|range(1, n + 1)| teaches inclusive versus exclusive bounds, the
single most common bug in transcribing mathematics into code, at the exact
moment it matters. That $S = k_B\log W$ needs $k_B$ and $W$ as arguments, unless
$k_B$ is a constant, teaches the difference between a variable and a parameter
and what ``in scope'' means. The self-hosting bootstrap is the capstone: a
student who reads the appendix sees the same program in two languages side by
side and can trace one into the other.

\paragraph{Checking (undergraduate and beyond).} Massot reports that students
using a proof assistant learn to write proofs that transfer to paper
\cite{massot2024}; Bottoni and colleagues report a positive influence on
students' understanding of what a proof is \cite{bottoni2025}. \texttt{lean4.py}
is smaller and earlier than either of those settings. It asks a student to
prove that equality is symmetric by writing one line of Python and one line of
\LaTeX{}, and then it tells them, precisely, whether they did. A student who
writes \verb|return x| to prove $x = x$ is told that they proved
$\mathrm{Nat}\to\mathrm{Nat}$ instead. That is a better tutor than most. And
because the kernel is readable, the student can then be asked to read it, and
to answer the question that separates people who have used a proof assistant
from people who understand one: what, exactly, is the machine trusting?

\paragraph{Computing.} The pair $m + 0 = m$ and $0 + n = n$ is the smallest
example we know of a lesson that only a proof checker can teach properly. Both
are obvious; one is proved by reduction and the other needs induction; and the
difference is a consequence of a choice made in the definition of addition,
which the student can change and watch the difficulty move to the other side.

\paragraph{Graduating.} The statement language uses Lean's spellings
deliberately---\verb|\forall|, braces for implicit arguments, \verb|refl|,
\verb|@| as \verb|explicit|---so that a student who moves to Lean and Mathlib
finds the same ideas with the same names. RosettaMath is a ramp, not a
destination. The destination is the system the FLT proof was checked with, and
the world Commelin et al.\ describe, in which mathematical training must
broaden to include the tools mathematics is now done with \cite{commelin2026}.

\subsection{Why local, open, and Python}

The design position of this project is that mathematics education software
should run on the student's own computer, should be open source, should be
written in the language the student is learning, and should be small enough to
modify. Each of these is a choice against a prevailing alternative, and we give
the reasons.

\paragraph{Against the client--server model.} Most educational software now
being built is a web application: a JavaScript front end, a server-side runtime,
an account, a network. That architecture was designed for software that needs
to be the same for millions of users and updated without their consent. It is a
poor fit for a classroom. It adds a server nobody in the room controls, a
network the school may not have, a login that expires, and two languages
neither the teacher nor the student is learning. A program that runs on a laptop
needs none of that. We have found that removing the server also removes most of
the complexity: the explorer's entire persistence layer is a dictionary.

\paragraph{Against the walled garden.} Proprietary mathematical software is
excellent and expensive, and its excellence is not the student's. A student
cannot read how it renders an equation, cannot correct an entry, cannot add
their teacher's notation, and takes nothing with them when the licence ends.
The alternative is not a worse proprietary tool but the same model that built
the systems mathematics itself now runs on: Linux, Python, \LaTeX{}, and now
Lean, whose stewards intend to become a foundation in the footsteps of Rust and
Linux \cite{simons2026}. These are also, not coincidentally, the tools of a
working scientist. A student who learns mathematics on them learns the
workbench along with the subject.

\paragraph{Against the phone.} We would like students to learn mathematics on
a machine with a keyboard, a file system, a terminal, and a compiler---a
machine on which the software they use is something they could, in principle,
have written. A tablet is a consumption device. The explorer does not run on
one, and this is a feature.

\paragraph{For Python.} The choice of Python is not a matter of taste. It is
the language secondary schools teach, the language undergraduate science
courses use, and the language in which the numerical libraries live. A
translator that emits Python emits something the student can already read.
A kernel written in Python is a kernel the student can already read. And
scipy is the reason the constants in an equation can be resolved to numbers by
name.

\paragraph{For smallness.} \texttt{rosettamath.py}, \texttt{rosettaui.py} and
\texttt{lean4.py} are 1,675, 3,171 and 1,945 lines respectively, in three
files, with one optional dependency (PyQt5) between them. This is small enough
that a teacher preparing a course on special relativity can add the Lorentz
transformations to the equation library in an hour, and small enough that a
student who wants the kernel to have inductive types can attempt it. Bastani's
trial \cite{bastani2025} and Wegerif's reading of it \cite{wegerif2026} suggest
that what a tool does to learning depends on how it is designed; a tool that the
learner can redesign is a tool whose design is itself a lesson.

\subsection{What this is not}

It is not a claim that RosettaMath replaces Lean, SymPy, Mathematica, or a
teacher. The subset is strict, the kernel is minimal, the knowledge base is a
few hundred entries, and the classifier is a signature matcher. It is a claim
about foundations: that the right base for mathematics education in the age of
machine proof is a set of small, local, open, readable tools in the student's
own language, and that these three are a working example.

\section{Beyond education: toward verified systems}
\label{sec:beyond}

The same small trusted kernel that makes \texttt{lean4.py} teachable makes it
a candidate for something else. The author's Crust project \cite{crust} is a
self-contained C, C++ and Rust toolchain that owns its whole stack---compiler,
emitted C, boot path, threading model---with no external compiler and, on
bare-metal targets, no operating system underneath. Its position is that every
credible safety story in systems programming, from JSF and MISRA to SPARK and
seL4 \cite{klein2009}, is a story of \emph{subtraction}: a smaller language made
load-bearing. Crust already carries contract-like declarations in its source,
such as assertions about which core a thread runs on, from which the compiler
derives register partitions and a context switcher; and its own paper on
proof-guided runtime checking \cite{crustsafety} argues for a toolchain small
enough to read.

The natural next step is to make the statement language of \texttt{lean4.py}
the contract language of Crust. A function contract would be a $\Pi$-type; a
call would be an application; the obligation to satisfy the contract would be
a theorem checked by the kernel, with the elaborator inferring the implicit
arguments as it does now. Crust's compiler is itself written in Python, so the
two projects share a host language and a philosophy. seL4 demonstrated that an
operating-system kernel can be proved correct; it did so with a large proof in
a large assistant. What a small assistant and a small toolchain make thinkable
is a verified system whose \emph{entire} trusted base---kernel, compiler and
checker---a single person can read. The Fermat formalisation was checked by a
second, independent kernel because the kernel is the only thing anyone must
believe \cite{anthropic2026flt}. That principle scales down as well as up.

This is future work, and we do not claim more than the shape of it here. But it
is why the kernel was built to be honest about what it trusts.

\section{Limitations}

The \LaTeX{} subset is deliberately small, and there are conventions it does
not read: summation and product notation are handled in pseudocode but not as
bare math-mode operators; matrices, integrals and limits are typeset but not
translated. The classifier is a signature matcher and will be confused by
equations that share symbols with unrelated ones; it says ``looks like'' for a
reason. The knowledge base is written by one author and has one author's
emphases. The kernel has inductive families with parameters and
indices, but no universe polymorphism, so every type lives at a fixed level and
each family needs a separate recursor for eliminating into $\mathrm{Type}$ and
into $\mathrm{Prop}$; its constraint solver is greedy across holes. The explorer's rendering is serviceable, not beautiful, and
degrades to a \texttt{pdflatex} image for anything demanding. None of these are
hidden in the code, and the test suites---the translator's bootstrap and 21
round-trip cases, the explorer's 53 checks and the kernel's 103---name what is
covered.

\section{Conclusion}

The equations published in a paper and the code executed in a simulation should
be the same artefact; a student should be able to click on any symbol in that
artefact and be told what it means; and a claim about the code should be
checkable by a kernel small enough to read. RosettaMath is a working example of
all three, built as three small, local, open Python programs, and offered as a
foundation for mathematics and programming education at every level in a
period when the proving of mathematics is being handed to machines. The
machines are learning to prove. The point of education is that people still
learn to understand, and understanding begins with being able to read the
notation, run it, and check it, on a machine of one's own.

%%APPENDIX%%

\small
\begin{thebibliography}{99}

\bibitem{knuth} Knuth, D.~E. (1984). Literate programming. \emph{The Computer
Journal}, 27(2), 97--111. \url{https://doi.org/10.1093/comjnl/27.2.97}

\bibitem{meurer} Meurer, A., Smith, C.~P., Paprocki, M., et al. (2017). SymPy:
symbolic computing in Python. \emph{PeerJ Computer Science}, 3, e103.
\url{https://doi.org/10.7717/peerj-cs.103}

\bibitem{poore} Poore, G.~M. (2015). PythonTeX: reproducible documents with
\LaTeX{}, Python, and more. \emph{Computational Science \& Discovery}, 8(1),
014010. \url{https://doi.org/10.1088/1749-4699/8/1/014010}

\bibitem{coquand1988} Coquand, T., \& Huet, G. (1988). The Calculus of
Constructions. \emph{Information and Computation}, 76(2--3), 95--120.

\bibitem{debruijn1972} de Bruijn, N.~G. (1972). Lambda calculus notation with
nameless dummies, a tool for automatic formula manipulation, with application
to the Church--Rosser theorem. \emph{Indagationes Mathematicae}, 34(5),
381--392.

\bibitem{miller1991} Miller, D. (1991). A logic programming language with
lambda-abstraction, function variables, and simple unification. \emph{Journal
of Logic and Computation}, 1(4), 497--536.

\bibitem{demoura2021} de Moura, L., \& Ullrich, S. (2021). The Lean 4 theorem
prover and programming language. In \emph{Automated Deduction -- CADE 28},
LNCS 12699, 625--635. Springer.

\bibitem{klein2009} Klein, G., Elphinstone, K., Heiser, G., et al. (2009).
seL4: formal verification of an OS kernel. In \emph{Proceedings of the 22nd
ACM Symposium on Operating Systems Principles (SOSP)}, 207--220.

\bibitem{massot2019} Massot, P. (2019). Lean in \LaTeX{}. \emph{The Xena
Project}, 11 February 2019.
\url{https://xenaproject.wordpress.com/2019/02/11/lean-in-latex/}

\bibitem{massot2024} Massot, P. (2024). Teaching mathematics using Lean and
controlled natural language. In \emph{15th International Conference on
Interactive Theorem Proving (ITP 2024)}, LIPIcs vol.~309, 27:1--27:19.
\url{https://doi.org/10.4230/LIPIcs.ITP.2024.27}

\bibitem{bottoni2025} Bottoni, M.~L., Cattaneo, A.~S., \& Sacikara, E. (2025).
Teaching ``Foundations of Mathematics'' with the LEAN theorem prover.
arXiv:2501.03352. \url{https://arxiv.org/abs/2501.03352}

\bibitem{wegerif2026} Wegerif, R. (2026). AI: researching the future of
education. \emph{Expanding Dialogic Space}, 17 August 2026.
\url{https://rupertwegerif.substack.com/p/ai-researching-the-future-of-education}

\bibitem{bastani2025} Bastani, H., Bastani, O., Sungu, A., Ge, H., Kabakc\i,
\"O., \& Mariman, R. (2025). Generative AI without guardrails can harm learning:
evidence from high school mathematics. \emph{Proceedings of the National
Academy of Sciences}, 122(26), e2422633122.
\url{https://doi.org/10.1073/pnas.2422633122}

\bibitem{commelin2026} Commelin, J., Jamnik, M., Ochigame, R., Taelman, L., \&
Venkatesh, A. (2026). Shaping the future of mathematics in the age of AI. To
appear in \emph{Notices of the American Mathematical Society}. arXiv:2603.24914.
\url{https://arxiv.org/abs/2603.24914}

\bibitem{letson2026} Letson, A., Sarra, L., Poiroux, A., et al. (2026). SorryDB:
can AI provers complete real-world Lean theorems? arXiv:2603.02668.
\url{https://arxiv.org/abs/2603.02668}

\bibitem{hubert2025} Hubert, T., Mehta, R., Sartran, L., et al. (2025).
Olympiad-level formal mathematical reasoning with reinforcement learning.
\emph{Nature}, 651, 607--613. \url{https://doi.org/10.1038/s41586-025-09833-y}

\bibitem{ren2025} Ren, Z.~Z., Shao, Z., Song, J., et al. (2025).
DeepSeek-Prover-V2: advancing formal mathematical reasoning via reinforcement
learning for subgoal decomposition. arXiv:2504.21801.
\url{https://arxiv.org/abs/2504.21801}

\bibitem{wang2025kimina} Wang, H., Unsal, M., Lin, X., et al. (2025).
Kimina-Prover Preview: towards large formal reasoning models with reinforcement
learning. arXiv:2504.11354. \url{https://arxiv.org/abs/2504.11354}

\bibitem{chen2025seed} Chen, L., Gu, J., Huang, L., et al. (2025). Seed-Prover:
deep and broad reasoning for automated theorem proving. arXiv:2507.23726.
\url{https://arxiv.org/abs/2507.23726}

\bibitem{lin2025goedel} Lin, Y., Tang, S., Lyu, B., et al. (2025).
Goedel-Prover-V2: scaling formal theorem proving with scaffolded data synthesis
and self-correction. arXiv:2508.03613. \url{https://arxiv.org/abs/2508.03613}

\bibitem{achim2025aristotle} Achim, T., Best, A., Bietti, A., et al. (2025).
Aristotle: IMO-level automated theorem proving. arXiv:2510.01346.
\url{https://arxiv.org/abs/2510.01346}

\bibitem{mathinc2025gauss} Math, Inc. (2025). Introducing Gauss, an agent for
autoformalization. September 2025. \url{https://www.math.inc/gauss}

\bibitem{anthropic2026flt} Anthropic (2026). Formalizing Fermat's Last Theorem.
September 2026. \url{https://www.anthropic.com/research/formalizing-fermats-last-theorem}

\bibitem{newscientist2026flt} \emph{New Scientist} (2026). Fermat's last
theorem formalised by AI agents in just 11 days.
\url{https://www.newscientist.com/article/2587839-fermats-last-theorem-formalised-by-ai-agents-in-just-11-days/}

\bibitem{buzzard2026counter} Buzzard, K. (2026). Human mathematicians are being
outcounterexampled. \emph{The Xena Project}, 20 July 2026.
\url{https://xenaproject.wordpress.com/2026/07/20/human-mathematicians-are-being-outcounterexampled/}

\bibitem{simons2026} Simons Foundation (2026). From trust to verification:
Lean's impact on mathematics. 23 June 2026.
\url{https://www.simonsfoundation.org/2026/06/23/from-trust-to-verification-leans-impact-on-mathematics/}

\bibitem{sloman2026} Sloman, L. (2026). In math, rigor is vital. But are
digitized proofs taking it too far? \emph{Quanta Magazine}, 25 March 2026.
\url{https://www.quantamagazine.org/in-math-rigor-is-vital-but-are-digitized-proofs-taking-it-too-far-20260325/}

\bibitem{peterman2024} Peterman, R. (2024). Creator of Lean: the end of
handwritten proofs and code. Interview with Leonardo de Moura, \emph{The
Developing Dev}.
\url{https://www.developing.dev/p/creator-of-lean-the-end-of-handwritten}

\bibitem{crust} Hartshorn, B.~S. (2026). Crust: a unified C/C++/Rust compiler
environment for systems programming. \url{https://github.com/brentharts/crust}

\bibitem{crustsafety} Hartshorn, B.~S. (2026). Memory safety where it is needed:
proof-guided runtime checking in a toolchain small enough to read. SSRN.
\url{https://dx.doi.org/10.2139/ssrn.7396160}

\end{thebibliography}

\end{document}
